\documentclass{ximera}

\begin{document}
	\author{Stitz-Zeager}
	\xmtitle{Exercises for Polar Coordinates}{}

\mfpicnumber{1} \opengraphsfile{ExercisesforPolarCoordinates} % mfpic settings added 

\begin{question}
In Exercises \ref{polarpointgraphfirst} - \ref{polarpointgraphlast}, plot the point given in polar coordinates and then give three different expressions for the point such that \hfill (a) $r < 0$ and $0 \leq \theta \leq 2\pi$, \hfill (b) $r > 0$ and $\theta \leq 0$ \hfill (c) $r > 0$ and $\theta \geq 2\pi$

\begin{problem}\label{polarpointgraphfirst}
$\left( 2, \frac{\pi}{3} \right)$    
\end{problem}

\begin{problem}
$\left( 5, \frac{7\pi}{4} \right)$
\end{problem}

\begin{problem}
$\left( \frac{1}{3}, \frac{3\pi}{2} \right)$ 
\end{problem}

\begin{problem}
$\left( \frac{5}{2}, \frac{5\pi}{6} \right)$   
\end{problem}

\begin{problem}
$\left( 12, -\frac{7\pi}{6} \right)$
\end{problem}

\begin{problem}
$\left( 3, -\frac{5\pi}{4} \right)$
\end{problem}

\begin{problem}
$\left( 2\sqrt{2}, -\pi \right)$ 
\end{problem}

\begin{problem}
$\left( \frac{7}{2}, -\frac{13\pi}{6} \right)$
\end{problem}

\begin{problem}
$\left( -20, 3\pi \right)$ 
\end{problem}

\begin{problem}
$\left( -4, \frac{5\pi}{4} \right)$
\end{problem}

\begin{problem}
$\left( -1, \frac{2\pi}{3} \right)$
\end{problem}

\begin{problem}
$\left( -3, \frac{\pi}{2} \right)$
\end{problem}

\begin{problem}
$\left( -3, -\frac{11\pi}{6} \right)$
\end{problem}

\begin{problem}
$\left( -2.5, -\frac{\pi}{4} \right)$
\end{problem}

\begin{problem}
$\left( -\sqrt{5}, -\frac{4\pi}{3} \right)$
\end{problem}

\begin{problem} \label{polarpointgraphlast}
$\left( -\pi, -\pi \right)$ 
\end{problem} 

\end{question}

\begin{question}
In Exercises \ref{polartorectfirst} - \ref{polartorectlast}, convert the point from polar coordinates into rectangular coordinates.  


\begin{problem}\label{polartorectfirst} 
$\left( 5, \frac{7\pi}{4} \right)$ 
\end{problem}

\begin{problem}
$\left( 2, \frac{\pi}{3} \right)$
\end{problem}

\begin{problem}
$\left( 11, -\frac{7\pi}{6} \right)$
\end{problem}

\begin{problem}
$\left( -20, 3\pi \right)$
\end{problem}

\begin{problem}
$\left( \frac{3}{5}, \frac{\pi}{2} \right)$
\end{problem}

\begin{problem}
$\left( -4, \frac{5\pi}{6} \right)$
\end{problem}

\begin{problem}
$\left( 9, \frac{7\pi}{2} \right)$
\end{problem}

\begin{problem}
$\left( -5, -\frac{9\pi}{4} \right)$
\end{problem}

\begin{problem}
$\left( 42, \frac{13\pi}{6} \right)$
\end{problem}

\begin{problem}
$\left( -117, 117\pi \right)$ 
\end{problem}

\begin{problem}
$\left( 6, \arctan(2) \right)$
\end{problem}

\begin{problem}
$\left(10, \arctan(3) \right)$ 
\end{problem}

\begin{problem}
$\left( -3, \arctan\left(\frac{4}{3}\right) \right)$
\end{problem}

\begin{problem}
$\left( 5, \arctan\left(-\frac{4}{3}\right) \right)$ 
\end{problem}

\begin{problem}
$\left( 2, \pi - \arctan\left(\frac{1}{2}\right)  \right)$ 
\end{problem}

\begin{problem}
$\left( -\frac{1}{2}, \pi - \arctan\left(5\right)  \right)$ 
\end{problem}

\begin{problem}
$\left( -1, \pi + \arctan\left(\frac{3}{4}\right) \right)$ 
\end{problem}

\begin{problem}
$\left( \frac{2}{3}, \pi + \arctan\left(2\sqrt{2}\right) \right)$
\end{problem}

\begin{problem}
 $\left( \pi, \arctan(\pi) \right)$
\end{problem}

\begin{problem} \label{polartorectlast}
$\left( 13, \arctan \left( \frac{12}{5} \right) \right)$ 
\end{problem}

\end{question}

\begin{question}
In Exercises \ref{recttopolarfirst} - \ref{recttopolarlast}, plot each point given in rectangular coordinates and convert to polar coordinates.  Choose $r \geq 0$ and $0 \leq \theta < 2\pi$.  

\begin{problem}\label{recttopolarfirst}
$(0, 5)$  
\end{problem}

\begin{problem}
$(3, \sqrt{3})$
\end{problem}

\begin{problem}
$(7, -7)$
\end{problem}

\begin{problem}
$(-3, -\sqrt{3})$
\end{problem}

\begin{problem}
$(-3, 0)$ 
\end{problem}

\begin{problem}
$\left( -\sqrt{2}, \sqrt{2} \right)$
\end{problem}

\begin{problem}
$\left( -4,-4\sqrt{3} \right)$ 
\end{problem}

\begin{problem}
$\left( \frac{\sqrt{3}}{4}, -\frac{1}{4} \right)$
\end{problem}

\begin{problem}
$\left( -\frac{3}{10}, -\frac{3\sqrt{3}}{10} \right)$
\end{problem}

\begin{problem}
$\left( -\sqrt{5}, -\sqrt{5} \right)$ 
\end{problem}

\begin{problem}
$(6,8)$
\end{problem}

\begin{problem}
$(\sqrt{5},2\sqrt{5})$
\end{problem}

\begin{problem}
$(-8,1)$ 
\end{problem}

\begin{problem}
$(-2\sqrt{10}, 6\sqrt{10})$ 
\end{problem}

\begin{problem}
$\left(-5, -12 \right)$
\end{problem}

\begin{problem}
$\left(-\frac{\sqrt{5}}{15}, -\frac{2\sqrt{5}}{15}  \right)$
\end{problem}

\begin{problem}
$\left(24, -7 \right)$ \vphantom{$\left(\frac{\sqrt{3}}{4}\right)$}
\end{problem}

\begin{problem}
$\left(12, -9\right)$
\end{problem}

\begin{problem}
$\left(\frac{\sqrt{2}}{4}, \frac{\sqrt{6}}{4}\right)$
\end{problem}

\begin{problem}\label{recttopolarlast}
$\left(-\frac{\sqrt{65}}{5}, \frac{2\sqrt{65}}{5}\right)$ 
\end{problem}

\end{question}


\begin{question}

In Exercises \ref{equrecttopolarfirst} - \ref{equrecttopolarlast}, convert the equation from rectangular coordinates into polar coordinates.  Solve for $r$ in all but \ref{solvethetaone} through \ref{solvethetafour}.  In Exercises \ref{solvethetaone} - \ref{solvethetafour},  solve for $\theta$

\begin{problem}\label{equrecttopolarfirst}
$x = 6$ 
\end{problem}

\begin{problem}
$x = -3$
\end{problem}

\begin{problem}
$y = 7$
\end{problem}

\begin{problem}\label{solvethetaone}
$y = 0$  
\end{problem}

\begin{problem}
$y = -x$
\end{problem}

\begin{problem}
$y = x\sqrt{3}$
\end{problem}

\begin{problem} \label{solvethetafour} 
$y = 2x$
\end{problem}

\begin{problem}
$x^2 + y^2 = 25$
\end{problem}

\begin{problem}
$x^{2} + y^{2} = 117$
\end{problem}

\begin{problem}
$y = 4x - 19$
\end{problem}

\begin{problem}
$x = 3y + 1$
\end{problem}

\begin{problem}
$y = -3x^{2}$
\end{problem}

\begin{problem}
$4x = y^2$
\end{problem}

\begin{problem}
$x^2 + y^2 - 2y = 0$
\end{problem}

\begin{problem}
$x^2 -4x + y^2 = 0$
\end{problem}

\begin{problem}
$x^2 + y^2 = x$
\end{problem}

\begin{problem}
$y^2 = 7y - x^2$
\end{problem}

\begin{problem}
$(x+2)^2 + y^2 = 4$
\end{problem}

\begin{problem}
$x^{2} + (y - 3)^{2} = 9$
\end{problem}

\begin{problem}\label{equrecttopolarlast}
$4x^2 + 4\left( y - \frac{1}{2} \right)^2 = 1$ 
\end{problem}

\end{question}

\begin{question}
In Exercises \ref{equpolartorectfirst} - \ref{equpolartorectlast}, convert the equation from polar coordinates into rectangular coordinates.  

\begin{problem} \label{equpolartorectfirst}
$r = 7$ 
\end{problem}

\begin{problem}
$r = -3$
\end{problem}

\begin{problem}
$r = \sqrt{2}$
\end{problem}

\begin{problem}
$\theta = \frac{\pi}{4}$
\end{problem}

\begin{problem}
$\theta = \frac{2\pi}{3}$
\end{problem}

\begin{problem}
$\theta = \pi$
\end{problem}

\begin{problem}
$\theta = \frac{3\pi}{2}$
\end{problem}

\begin{problem}
$r = 4\cos(\theta)$
\end{problem}

\begin{problem}
$5r = \cos(\theta)$
\end{problem}

\begin{problem}
$r = 3\sin(\theta)$
\end{problem}

\begin{problem}
$r = -2\sin(\theta)$
\end{problem}

\begin{problem}
$r = 7\sec(\theta)$
\end{problem}

\begin{problem}
$12r = \csc(\theta)$
\end{problem}

\begin{problem}
$r = -2\sec(\theta)$
\end{problem}

\begin{problem}
$r = -\sqrt{5} \csc(\theta)$
\end{problem}

\begin{problem}
$r = 2\sec(\theta)\tan(\theta)$
\end{problem}

\begin{problem}
$r = -\csc(\theta) \cot(\theta)$
\end{problem}

\begin{problem}
$r^{2} = \sin(2\theta)$
\end{problem}

\begin{problem}
$r = 1 - 2\cos(\theta)$
\end{problem}

\begin{problem} \label{equpolartorectlast}
$r = 1 + \sin(\theta)$ 
\end{problem}

\end{question}

\begin{problem} Convert the origin $(0,0)$ into polar coordinates in four different ways.
\end{problem}

\begin{problem}
With the help of your classmates, use the Law of Cosines to develop a formula for the distance between two points in polar coordinates.
\end{problem}

\end{document}
